\chapter{Michael H. Freedman --- Fields Medal 1986}\label{ch1986michaelfreedman:michael-h.-freedman-fields-medal-1986-a-comprehensive-report}

\begin{center}\rule{0.5\linewidth}{0.5pt}\end{center}

\textbf{Prepared:} July 2026 \textbf{Subject:} Michael Hartley Freedman
(b. April 21, 1951, Los Angeles, California, USA) \textbf{Position:}
Professor, Microsoft Station Q, University of California, Santa Barbara
(since 2005); Professor Emeritus, UC San Diego \textbf{Doctoral
Advisor:} William Browder (Princeton University) \textbf{Fields Medal
Citation:} \emph{``For his work on the topological classification of
four-manifolds, including the proof of the Poincare conjecture for
four-dimensional manifolds (the topological case) and the classification
of simply connected four-manifolds.''}

\begin{center}\rule{0.5\linewidth}{0.5pt}\end{center}

\section{1. Introduction}\label{ch1986michaelfreedman:introduction}

Michael Freedman was awarded the Fields Medal in 1986 for his
substantial work on the topology of four-manifolds. He proved the
\textbf{topological Poincare conjecture in dimension 4} --- that every
homotopy 4-sphere is homeomorphic to \(S^4\) --- and gave a complete
\textbf{classification of simply connected topological four-manifolds}
up to homeomorphism.

\begin{prereqs} Algebraic topology (intersection forms, characteristic classes), basic differential topology (handle decompositions). \end{prereqs}

\begin{objectives} State the topological classification theorem for simply connected 4-manifolds; explain the role of intersection forms; describe the relation to Donaldson's theorem. \end{objectives}


Freedman's work showed that the topology of four-manifolds is
substantially different from that of higher-dimensional manifolds. While
the smooth structure of four-manifolds is wild and exotic (as
demonstrated by Donaldson's work, which earned him a share of the same
Fields Medal), the topological classification is surprisingly clean and
complete. This divergence between the topological and smooth categories
is unique to dimension 4 --- in all other dimensions, these two
categories are much more closely aligned.

The 1986 Fields Medal was awarded to three mathematicians --- Freedman,
Donaldson, and Faltings --- whose work together transformed algebraic
and differential topology. Freedman's topological classification and
Donaldson's smooth obstructions are complementary: Freedman shows what
is possible topologically, and Donaldson shows what is impossible
smoothly. Together, they revealed the notable richness of
four-dimensional geometry.

In his later career, Freedman made a notable transition from pure
mathematics to applied mathematics and physics. He founded Microsoft
Station Q at UC Santa Barbara and became a pioneer in
\textbf{topological quantum computing}, using the mathematical
structures he discovered in four-manifold topology --- modular tensor
categories, anyon models, and link invariants --- as the basis for a
substantially new approach to quantum computation.

\begin{center}\rule{0.5\linewidth}{0.5pt}\end{center}

\section{2. Biographical Background}\label{ch1986michaelfreedman:biographical-background}

\subsection{2.1 Early Life and
Education}\label{ch1986michaelfreedman:early-life-and-education}

Michael Freedman was born in Los Angeles in 1951 into a creative family
--- his father was a fashion designer and his mother a painter. Despite
this artistic background, he was drawn to mathematics at an early age.
He studied at the University of California, Santa Cruz as an
undergraduate, where his mathematical talents became evident, before
moving to Princeton University for doctoral work.

\subsection{2.2 Doctoral Work}\label{ch1986michaelfreedman:doctoral-work}

At Princeton, Freedman studied under William Browder, one of the leading
figures in algebraic and differential topology. His PhD dissertation,
completed in 1973, dealt with problems in four-manifold topology,
foreshadowing his later significant contributions.

\subsection{2.3 Academic Career}\label{ch1986michaelfreedman:academic-career}

\begin{itemize}
\tightlist
\item
 \textbf{1973:} PhD, Princeton University, under William Browder.
\item
 \textbf{1973--1975:} Assistant Professor, UC Berkeley.
\item
 \textbf{1975--1995:} Professor, UC San Diego, where he produced his
 Fields Medal work on the topological classification of four-manifolds.
\item
 \textbf{1982:} Proof of the topological Poincare conjecture in
 dimension 4.
\item
 \textbf{1984:} Publication of the complete classification of simply
 connected topological four-manifolds.
\item
 \textbf{1986:} Fields Medal awarded at the ICM in Berkeley.
\item
 \textbf{1995--2005:} Professor, University of Washington, Seattle.
\item
 \textbf{1997:} Elected to the National Academy of Sciences.
\item
 \textbf{2003:} Elected to the American Academy of Arts and Sciences.
\item
 \textbf{2005--present:} Microsoft Research, Station Q, UC Santa
 Barbara.
\item
 \textbf{2006:} National Medal of Science.
\item
 \textbf{Additional:} Visited the Institute for Advanced Study, the
 Mathematical Sciences Research Institute, and numerous other
 institutions throughout his career.
\end{itemize}

\begin{center}\rule{0.5\linewidth}{0.5pt}\end{center}

\section{3. The Poincare Conjecture}\label{ch1986michaelfreedman:the-poincare-conjecture}

\subsection{3.1 The Statement}\label{ch1986michaelfreedman:the-statement}

The \textbf{Poincare conjecture} states that every simply connected
closed three-manifold is homeomorphic to the 3-sphere \(S^3\). It was
formulated by Henri Poincare in 1904 as a fundamental question in
topology. The 3-dimensional case was proved by Grigori Perelman in 2003,
earning him the Fields Medal in 2006. However, there is a natural
higher-dimensional generalization of the conjecture that was the subject
of intense study throughout the 20th century.

\subsection{3.2 The Generalized Version}\label{ch1986michaelfreedman:the-generalized-version}

The \textbf{generalized Poincare conjecture} asks whether every homotopy
\(n\)-sphere --- a closed \(n\)-manifold with the homotopy type of
\(S^n\) --- is homeomorphic to \(S^n\). The history of its solution is
one of the great sagas of 20th-century mathematics: - For \(n \geq 5\),
proved by Stephen Smale in the 1960s (Fields Medal 1966) using
handlebody theory and the h-cobordism theorem. - For \(n = 2\), known
since the classification of surfaces in the 19th century. - For
\(n = 3\), proved by Perelman (Fields Medal 2006). - For \(n = 4\),
proved by Freedman.

The dimension 4 case was the last to fall and required completely new
techniques. The Whitney trick --- which Smale used so effectively in
dimensions \(\geq 5\) to eliminate intersections --- fails in dimension
4 for smooth embeddings. This failure is the reason dimension 4 is
special.

\subsection{3.3 Freedman's Proof}\label{ch1986michaelfreedman:freedmans-proof}

Freedman's proof of the 4-dimensional case uses three main innovations:
- \textbf{Casson handles:} A geometric construction that produces
topological (but not necessarily smooth) 2-handles. These are the
4-dimensional analogue of the handles Smale used in higher dimensions. -
\textbf{Controlled surgery theory:} A version of Wall's surgery theory
adapted to the special problems of dimension 4, where the Whitney trick
fails. - \textbf{The disk theorem:} A proof that Casson handles are
homeomorphic to standard handles, which is the key technical result.

\subsection{3.4 Significance}\label{ch1986michaelfreedman:significance}

The proof was an unexpected result. Many
topologists had believed the conjecture would be false in dimension 4,
and were astonished when Freedman proved it true. Importantly, the
smooth version of the Poincare conjecture in dimension 4 remains open to
this day --- we do not know whether a smooth homotopy 4-sphere is
diffeomorphic to \(S^4\).

\begin{center}\rule{0.5\linewidth}{0.5pt}\end{center}

\section{4. The Classification
Theorem}\label{ch1986michaelfreedman:the-classification-theorem}

\subsection{4.1 The Complete Set of
Invariants}\label{ch1986michaelfreedman:the-complete-set-of-invariants}

Freedman classified all simply connected topological four-manifolds
using just two invariants: - The \textbf{intersection form}
\(Q_X: H_2(X; \mathbb{Z}) \times H_2(X; \mathbb{Z}) \to \mathbb{Z}\),
which is a unimodular symmetric bilinear form. This encodes the
algebraic intersection numbers of 2-dimensional cycles in the manifold.
The form is symmetric because \(Q_X(\alpha, \beta)\) counts signed
intersections of 2-cycles, and it is unimodular by Poincare duality. -
The \textbf{Kirby--Siebenmann invariant}
\(\kappa(X) \in \mathbb{Z}/2\mathbb{Z}\), which is the primary
obstruction to the existence of a smooth structure on \(X\).

\subsection{4.2 The Full Statement}\label{ch1986michaelfreedman:the-full-statement}

Freedman's classification theorem is remarkably simple and complete. It
states: for every unimodular symmetric bilinear form \(Q\) over
\(\mathbb{Z}\), there exists a simply connected topological
four-manifold \(X\) with intersection form isomorphic to \(Q\) if and
only if either \(Q\) is even or \(\kappa(X) = 0\). Moreover, the
classification is unique in the following sense: - When \(Q\) is even,
there are exactly two distinct homeomorphism types --- one with
\(\kappa = 0\) (smoothable in principle) and one with \(\kappa = 1\)
(not smoothable). - When \(Q\) is odd, there is exactly one
homeomorphism type, with \(\kappa\) determined by \(Q\).

\subsection{4.3 Consequences}\label{ch1986michaelfreedman:consequences}

The classification has significant implications for our understanding of
four-dimensional geometry: - Every simply connected topological
four-manifold is completely determined by its intersection form and the
Kirby--Siebenmann invariant --- these two invariants capture all
topological information. - Algebraic intersection forms that are not
realized by any smooth four-manifold (such as the \(E_8\) form) are
realized topologically, highlighting the gap between the topological and
smooth categories. - There are vastly more topological four-manifolds
than smooth ones --- most topological four-manifolds do not admit any
smooth structure. - The classification is complete: there are no further
invariants beyond \(Q_X\) and \(\kappa(X)\), which is a striking
simplification.

\subsection{4.4 Comparison with Higher
Dimensions}\label{ch1986michaelfreedman:comparison-with-higher-dimensions}

In dimensions \(\geq 5\), the classification of simply connected
manifolds is much more complicated, involving homotopy theory,
characteristic classes, and surgery theory. The simplicity of Freedman's
classification in dimension 4 is surprising and beautiful, even though
it applies only to the topological category.

\begin{center}\rule{0.5\linewidth}{0.5pt}\end{center}

\section{5. The Intersection Form}\label{ch1986michaelfreedman:the-intersection-form}

\subsection{5.1 Definition}\label{ch1986michaelfreedman:definition}

For a closed oriented four-manifold \(X\), the \textbf{intersection
form} is the symmetric bilinear pairing on middle homology:

\[
Q_X(\alpha, \beta) = \langle \alpha \cup \beta, [X] \rangle
\]

\subsection{5.2 Properties}\label{ch1986michaelfreedman:properties}

\begin{itemize}
\tightlist
\item
 Symmetric: \(Q_X(\alpha, \beta) = Q_X(\beta, \alpha)\).
\item
 Unimodular: \(\det(Q_X) = \pm 1\) by Poincare duality over
 \(\mathbb{Z}\).
\item
 Even/odd: even if all self-intersection numbers are even, odd
 otherwise.
\item
 The signature \(\sigma(X) = b_2^+(X) - b_2^-(X)\) is the difference
 between the numbers of positive and negative eigenvalues of \(Q_X\).
\item
 The form determines the Euler characteristic:
 \(\chi(X) = 2 + b_2(X)\).
\end{itemize}

\subsection{5.3 Key Examples}\label{ch1986michaelfreedman:key-examples}

\begin{itemize}
\tightlist
\item
 \(S^4\): trivial form of rank 0 (the 4-sphere has no
 middle-dimensional homology).
\item
 \(\mathbb{CP}^2\): the positive definite form {[}1{]} (rank 1,
 signature 1).
\item
 \(\overline{\mathbb{CP}}^2\): the negative definite form {[}-1{]}
 (rank 1, signature -1).
\item
 \(S^2 \times S^2\): the hyperbolic form
 \(H = \begin{pmatrix} 0 & 1 \\ 1 & 0 \end{pmatrix}\) (rank 2,
 signature 0).
\item
 K3 surface: \(2E_8 \oplus 3H\) (rank 22, signature -16) --- the most
 important example of a simply connected 4-manifold with an even
 intersection form.
\end{itemize}

\subsection{5.4 Classification of Forms}\label{ch1986michaelfreedman:classification-of-forms}

The classification of unimodular symmetric bilinear forms over
\(\mathbb{Z}\) is a classical result: - \textbf{Indefinite forms:}
Classified by rank, signature, and parity. - \textbf{Definite forms:}
More complicated --- there are many definite forms, and they are still
not fully classified. - \textbf{The \(E_8\) form:} The smallest even
positive definite form, of rank 8 and signature 8.

\begin{center}\rule{0.5\linewidth}{0.5pt}\end{center}

\section{6. The Kirby--Siebenmann
Invariant}\label{ch1986michaelfreedman:the-kirbysiebenmann-invariant}

\subsection{6.1 Definition}\label{ch1986michaelfreedman:definition-1}

The \textbf{Kirby--Siebenmann invariant}
\(\kappa(X) \in \mathbb{Z}/2\mathbb{Z}\) is the primary obstruction to
the existence of a smooth structure on a topological four-manifold
\(X\). It was developed by Robion Kirby and Larry Siebenmann in their
seminal work on the triangulation of manifolds in the 1970s.

\subsection{6.2 Properties}\label{ch1986michaelfreedman:properties-1}

\begin{itemize}
\tightlist
\item
 \(\kappa(X) = 0\) if and only if \(X \times \mathbb{R}\) is smoothable
 (i.e., the product with \(\mathbb{R}\) admits a smooth structure).
\item
 The invariant is additive under connected sum:
 \(\kappa(X \# Y) = \kappa(X) + \kappa(Y) \pmod{2}\).
\item
 For any smoothable manifold, \(\kappa(X) = 0\).
\item
 For even intersection forms, \(\kappa(X) \equiv \sigma(X)/8 \pmod{2}\)
 where \(\sigma(X)\) is the signature.
\item
 The invariant is a homeomorphism invariant: if \(X \cong Y\)
 topologically, then \(\kappa(X) = \kappa(Y)\).
\end{itemize}

\subsection{6.3 Role in Classification}\label{ch1986michaelfreedman:role-in-classification}

The Kirby--Siebenmann invariant interacts with the intersection form in
the classification: - For odd intersection forms, \(\kappa\) is
determined by \(Q\) (specifically, \(\kappa(X) = 0\) when the signature
is divisible by 8 and \(\kappa(X) = 1\) otherwise). - For even
intersection forms, both values \(\kappa = 0\) and \(\kappa = 1\) occur,
yielding two distinct homeomorphism types with the same intersection
form. One of these is smoothable (\(\kappa = 0\)) and the other is not
(\(\kappa = 1\)).

\subsection{6.4 Historical Context}\label{ch1986michaelfreedman:historical-context}

The Kirby--Siebenmann invariant was a major development in topology. It
completed the understanding of the relationship between topological and
smooth manifolds and showed that the obstruction to smoothing is
codified in a single \(\mathbb{Z}/2\) class.

\begin{center}\rule{0.5\linewidth}{0.5pt}\end{center}

\section{7. Casson Handles}\label{ch1986michaelfreedman:casson-handles}

\subsection{7.1 Origin}\label{ch1986michaelfreedman:origin}

\textbf{Casson handles} were introduced by Andrew Casson in the 1970s as
a way to construct a 2-handle in a four-manifold without using the
Whitney trick, which is unavailable in dimension 4 for smooth
embeddings. Casson's insight was that while the Whitney disk cannot be
made smooth and embedded, it can be made topologically embedded through
an infinite construction.

\subsection{7.2 Construction}\label{ch1986michaelfreedman:construction}

A Casson handle is built by an infinite process. Start with a 2-handle
\(D^2 \times D^2\) attached to a four-manifold. The attaching region for
the next layer is the boundary of this handle:
\(\partial D^2 \times D^2 \cong S^1 \times D^2\). Inside this solid
torus, one embeds a framed link --- a collection of circles linked in
the pattern of the Borromean rings --- and attaches more 2-handles to
these circles. This process is repeated infinitely, producing an
infinite tower of handles, each layer attached to the boundary of the
previous layer.

\subsection{7.3 Properties}\label{ch1986michaelfreedman:properties-2}

\begin{itemize}
\tightlist
\item
 A Casson handle is homeomorphic to the standard open 2-handle
 \(D^2 \times \mathbb{R}^2\). This is Freedman's disk theorem.
\item
 It is not necessarily diffeomorphic to the standard handle, and in
 general it is not.
\item
 The boundary of a Casson handle is a topological 3-sphere but not
 necessarily a smooth 3-sphere.
\item
 The Casson handle encodes a deep relationship between the topology of
 4-manifolds and the theory of knots and links in 3-manifolds.
\end{itemize}

\subsection{7.4 Freedman's Contribution}\label{ch1986michaelfreedman:freedmans-contribution}

Freedman proved that Casson handles are homeomorphic to standard
handles, completing Casson's program. This proof is one of the deepest
and most intricate results in topology. It involves a careful analysis
of the infinite tree structure of the Casson handle and a sophisticated
application of Bing's shrinking technique.

\subsection{7.5 Why Casson Handles Work}\label{ch1986michaelfreedman:why-casson-handles-work}

Casson handles work because the infinite process of attaching handles
``spreads out'' the singularities that prevent the Whitney trick from
working smoothly. The end result is a topological embedding that cannot
be smoothed, reflecting the fundamental difference between the
topological and smooth categories in dimension 4.

\begin{center}\rule{0.5\linewidth}{0.5pt}\end{center}

\section{8. The Disk Theorem}\label{ch1986michaelfreedman:the-disk-theorem}

\subsection{8.1 Statement}\label{ch1986michaelfreedman:statement}

The \textbf{disk theorem} (also called the Freedman embedding theorem)
is the central technical result of Freedman's work. It states that every
Casson handle is homeomorphic to the standard open 2-handle
\(D^2 \times \mathbb{R}^2\). Consequently, any topologically embedded
2-sphere in a four-manifold can be thickened to a topological 2-handle.

\subsection{8.2 Significance}\label{ch1986michaelfreedman:significance-1}

The disk theorem is the engine that drives the entire classification. It
provides: - A method for constructing embedded disks in four-manifolds,
bypassing the failure of the Whitney trick. - A topological handlebody
theory for dimension 4, analogous to Smale's handlebody theory in higher
dimensions. - The ability to perform surgery in the topological category
of dimension 4. - The bridge between Casson's geometric ideas and the
complete classification theorem.

\subsection{8.3 Proof Outline}\label{ch1986michaelfreedman:proof-outline}

The proof of the disk theorem involves several deep steps: 1. Analyzing
the infinite tree structure of the Casson handle --- each node of the
tree corresponds to a layer of the handle. 2. Using Bing's shrinking
technique to collapse the handle to a standard one by ``shrinking'' the
non-standard parts to points. 3. Controlling the amount of shrinking at
each level of the infinite construction to ensure the final result is a
homeomorphism. 4. Verifying that the shrinking process does not create
new pathological behavior.

\subsection{8.4 Impact}\label{ch1986michaelfreedman:impact}

The disk theorem is one of the landmark results in topology. It shows
that the wild, infinite topology of Casson handles can be tamed, and it
provides the key that unlocks the full classification of four-manifolds.

\begin{center}\rule{0.5\linewidth}{0.5pt}\end{center}

\section{9. Controlled Surgery Theory}\label{ch1986michaelfreedman:controlled-surgery-theory}

\subsection{9.1 Classical Surgery}\label{ch1986michaelfreedman:classical-surgery}

\textbf{Surgery theory} (developed by Browder, Novikov, Sullivan, and
Wall in the 1960s) is a powerful method for constructing and classifying
manifolds. The basic operation is surgery: remove a thickened sphere
\(S^k \times D^{n-k}\) from an \(n\)-manifold and glue in
\(D^{k+1} \times S^{n-k-1}\) along the boundary. This changes the
manifold in a controlled way. The theory works perfectly in dimensions
\(\geq 5\) because the Whitney trick can eliminate unwanted
intersections between embedded spheres.

\subsection{9.2 The Obstacle in Dimension
4}\label{ch1986michaelfreedman:the-obstacle-in-dimension-4}

The Whitney trick fails in dimension 4 for smooth embeddings. Two
2-dimensional submanifolds in a 4-manifold generically intersect in
isolated points. In higher dimensions, these intersections can be
removed by the Whitney trick: one finds an embedded 2-disk (the Whitney
disk) connecting the intersection points and pushes one submanifold
across this disk. In dimension 4, the Whitney disk is 2-dimensional, and
it generically intersects the submanifolds it is meant to separate,
creating new problems.

\subsection{9.3 Freedman's Solution}\label{ch1986michaelfreedman:freedmans-solution}

Freedman developed a \textbf{controlled surgery theory} that overcomes
the failure of the Whitney trick: - Uses topological (Casson handle)
embeddings instead of smooth ones --- the Casson handle provides a
topological Whitney disk. - Works in a controlled, small-scale setting
where the geometry can be carefully managed. - Allows the classification
to proceed despite the failure of the smooth Whitney trick. - This is
why Freedman's classification is topological rather than smooth.

\subsection{9.4 Consequences}\label{ch1986michaelfreedman:consequences-1}

The controlled surgery theory enables: - The construction of
four-manifolds with prescribed intersection forms. - The classification
of simply connected four-manifolds. - The topological h-cobordism
theorem in dimension 4.

\begin{center}\rule{0.5\linewidth}{0.5pt}\end{center}

\section{10. The h-Cobordism Theorem in Dimension
4}\label{ch1986michaelfreedman:the-h-cobordism-theorem-in-dimension-4}

\subsection{10.1 The Classical Result}\label{ch1986michaelfreedman:the-classical-result}

The \textbf{h-cobordism theorem} (Smale, 1960s) is one of the landmark
results in differential topology. It states: if \(W\) is an h-cobordism
between two simply connected manifolds \(M\) and \(N\) of dimension
\(\geq 5\), then \(W\) is diffeomorphic to the product
\(M \times [0, 1]\). This theorem implies the Poincare conjecture in
dimensions \(\geq 5\) and is the foundation of surgery theory.

\subsection{10.2 The Topological Version in Dimension
4}\label{ch1986michaelfreedman:the-topological-version-in-dimension-4}

Freedman proved the topological version of the h-cobordism theorem in
dimension 4: if \(W\) is a topological h-cobordism between simply
connected four-manifolds, then \(W\) is homeomorphic to a product. The
proof uses Casson handles to replace the smooth handles that would be
used in higher dimensions.

\subsection{10.3 The Smooth Version Fails in Dimension
4}\label{ch1986michaelfreedman:the-smooth-version-fails-in-dimension-4}

The smooth version of the h-cobordism theorem \textbf{fails} in
dimension 4. Donaldson's gauge-theoretic work provides explicit
counterexamples: there exist smooth h-cobordisms between smooth
four-manifolds that are not diffeomorphic to a product. The reason is
that Casson handles are only topological, not smooth --- the Whitney
disk can be constructed topologically but not smoothly.

\subsection{10.4 Philosophical
Implications}\label{ch1986michaelfreedman:philosophical-implications}

The failure of the smooth h-cobordism theorem in dimension 4 shows that
smooth four-manifolds are substantially different from smooth manifolds
in other dimensions. This makes dimension 4 both frustrating and
fascinating: it is the dimension where the most exotic phenomena occur,
and it is the dimension we live in physically.

\subsection{10.5 Consequences}\label{ch1986michaelfreedman:consequences-2}

The failure of the smooth h-cobordism theorem has many consequences: -
Existence of exotic \(\mathbb{R}^4\)'s (smooth manifolds homeomorphic
but not diffeomorphic to \(\mathbb{R}^4\)). - Multiple smooth structures
on the same topological four-manifold. - The richness of gauge-theoretic
invariants in dimension 4.

\begin{center}\rule{0.5\linewidth}{0.5pt}\end{center}

\section{\texorpdfstring{11. The \(E_8\)
Manifold}{11. The E\_8 Manifold}}\label{ch1986michaelfreedman:the-e_8-manifold}

\subsection{11.1 The Form}\label{ch1986michaelfreedman:the-form}

The \(E_8\) form is the positive definite unimodular even symmetric
bilinear form of rank 8. It corresponds to the Cartan matrix of the
exceptional Lie algebra \(E_8\) and has signature 8. It is the smallest
even positive definite form that is not a sum of hyperbolic forms.

\subsection{11.2 Topological Existence}\label{ch1986michaelfreedman:topological-existence}

Freedman showed that there exists a topological four-manifold with
intersection form \(E_8\). This manifold, called the \(E_8\) manifold,
is constructed by plumbing disk bundles over spheres according to the
\(E_8\) Dynkin diagram, then adding a Casson handle to make it simply
connected. The Casson handle is needed to kill the fundamental group of
the plumbed manifold without changing the intersection form.

\subsection{11.3 Non-Smoothability}\label{ch1986michaelfreedman:non-smoothability}

Donaldson proved that the \(E_8\) manifold does not admit any smooth
structure. This was one of the first and most striking applications of
Donaldson's gauge-theoretic invariants. It means that the intersection
form \(E_8\) --- which is perfectly valid as an algebraic form --- can
be realized topologically but not smoothly.

\subsection{11.4 Significance}\label{ch1986michaelfreedman:significance-2}

The \(E_8\) manifold plays a central role in the construction of exotic
\(\mathbb{R}^4\)'s. By removing a certain Casson handle from the \(E_8\)
manifold, one obtains an exotic copy of \(\mathbb{R}^4\). The \(E_8\)
manifold is also the simplest example of a topological four-manifold
with no smooth structure, and it illustrates the fundamental gap between
the topological and smooth categories in dimension 4.

\subsection{\texorpdfstring{11.5 The \(E_8\) Plumbed
Manifold}{11.5 The E\_8 Plumbed Manifold}}\label{ch1986michaelfreedman:the-e_8-plumbed-manifold}

The \(E_8\) manifold can be constructed by plumbing 8 copies of the
\(D^2\)-bundle over \(S^2\) according to the \(E_8\) Dynkin diagram. The
resulting manifold has boundary equal to the Poincare homology sphere
\(\Sigma(2,3,5)\), and it has the \(E_8\) intersection form. To make it
simply connected, one caps off the boundary with a Casson handle.

\begin{center}\rule{0.5\linewidth}{0.5pt}\end{center}

\section{\texorpdfstring{12. Exotic
\(\mathbb{R}^4\)}{12. Exotic \textbackslash mathbb\{R\}\^{}4}}\label{ch1986michaelfreedman:exotic-mathbbr4}

\subsection{12.1 Definition}\label{ch1986michaelfreedman:definition-2}

An \textbf{exotic \(\mathbb{R}^4\)} is a smooth manifold that is
homeomorphic to \(\mathbb{R}^4\) but not diffeomorphic to it. The
existence of such a manifold was one of the most surprising consequences
of the work of Freedman and Donaldson.

\subsection{12.2 Construction}\label{ch1986michaelfreedman:construction-1}

Combining Freedman's topological methods with Donaldson's gauge theory
yields exotic \(\mathbb{R}^4\)'s. The construction removes a Casson
handle from the \(E_8\) manifold, producing an open subset that is
homeomorphic to \(\mathbb{R}^4\). However, because the Casson handle is
not smoothly standard, this \(\mathbb{R}^4\) is not diffeomorphic to the
standard one.

\subsection{12.3 Uniqueness to Dimension
4}\label{ch1986michaelfreedman:uniqueness-to-dimension-4}

Dimension 4 is the only dimension in which \(\mathbb{R}^n\) admits
exotic smooth structures. In dimensions \(n \neq 4\), \(\mathbb{R}^n\)
has a unique smooth structure up to diffeomorphism. This makes dimension
4 truly special in differential topology.

\subsection{12.4 An Infinitude of
Exotics}\label{ch1986michaelfreedman:an-infinitude-of-exotics}

In fact, there are uncountably many distinct exotic \(\mathbb{R}^4\)'s.
This was shown by Taubes and others using further gauge-theoretic
methods. The family of exotic \(\mathbb{R}^4\)'s forms a rich continuum
of distinct smooth structures, all on the same underlying topological
space.

\subsection{\texorpdfstring{12.5 Classification of Exotic
\(\mathbb{R}^4\)'s}{12.5 Classification of Exotic \textbackslash mathbb\{R\}\^{}4's}}\label{ch1986michaelfreedman:classification-of-exotic-mathbbr4s}

There are two types of exotic \(\mathbb{R}^4\)'s: - \textbf{Small exotic
\(\mathbb{R}^4\):} These embed in the standard \(\mathbb{R}^4\) and are
obtained by Casson handle constructions. - \textbf{Large exotic
\(\mathbb{R}^4\):} These do not embed in the standard \(\mathbb{R}^4\)
and are obtained from non-smoothable four-manifolds via more elaborate
constructions.

\subsection{12.6 The Enduring Mystery}\label{ch1986michaelfreedman:the-enduring-mystery}

Exotic \(\mathbb{R}^4\) remains one of the most surprising phenomena in
mathematics. It shows that the smooth structure of the simplest
non-compact four-manifold is not unique, and it continues to be an
active area of research.

\begin{center}\rule{0.5\linewidth}{0.5pt}\end{center}

\section{13. The Borromean Rings and Casson
Handles}\label{ch1986michaelfreedman:the-borromean-rings-and-casson-handles}

\subsection{13.1 The Borromean Rings}\label{ch1986michaelfreedman:the-borromean-rings}

The \textbf{Borromean rings} are three linked circles in \(S^3\) that
are pairwise unlinked (any two can be separated) but the three together
cannot be separated. They are the simplest Brunnian link and have been
used as a symbol of interdependence for centuries, appearing in
heraldry, religious art, and now mathematics.

\subsection{13.2 Role in Casson Handles}\label{ch1986michaelfreedman:role-in-casson-handles}

The infinite construction of a Casson handle is controlled by a tree of
Borromean ring configurations. At each stage of the construction, the
attaching circles of the next layer of handles are linked in the pattern
of the Borromean rings. The combinatorics of this tree of Borromean
rings determines whether the Casson handle is homeomorphic to the
standard handle.

\subsection{13.3 Why the Borromean Rings
Appear}\label{ch1986michaelfreedman:why-the-borromean-rings-appear}

The appearance of the Borromean rings is not accidental. In a
4-manifold, the attaching circle of a 2-handle is a circle in the
3-dimensional boundary. When two handles are attached, their attaching
circles can be linked. The Borromean configuration is the simplest
non-trivial linking pattern that allows the infinite construction to
proceed without creating obstructions.

\subsection{13.4 The Tree Structure}\label{ch1986michaelfreedman:the-tree-structure}

The Borromean rings at each stage of the Casson handle form nodes in an
infinite binary tree. The branching structure of this tree encodes the
topological complexity of the Casson handle. If the tree has a certain
``regular'' structure, the Casson handle is standard; otherwise, it may
be exotic.

\begin{center}\rule{0.5\linewidth}{0.5pt}\end{center}

\section{14. The Slice-Ribbon Problem}\label{ch1986michaelfreedman:the-slice-ribbon-problem}

\subsection{14.1 Slice Knots}\label{ch1986michaelfreedman:slice-knots}

A knot in \(S^3\) is \textbf{slice} if it bounds a smoothly embedded
disk in the 4-ball \(D^4\). Determining which knots are slice is a
central problem in knot theory with connections to four-manifold
topology and algebraic geometry.

\subsection{14.2 Topologically Slice
Knots}\label{ch1986michaelfreedman:topologically-slice-knots}

Freedman's disk theorem shows that many knots are \textbf{topologically
slice}: they bound a topologically embedded disk in \(D^4\). In
particular, every knot with Alexander polynomial 1 is topologically
slice. This includes many knots that were not known to be smoothly
slice.

\subsection{14.3 The Gap Between Topological and
Smooth}\label{ch1986michaelfreedman:the-gap-between-topological-and-smooth}

The gap between topological and smooth sliceness is a major theme in
knot theory. Many knots that are topologically slice are not smoothly
slice, and detecting this difference requires deep gauge-theoretic or
Floer-homological invariants such as the Ozsvath--Szabo knot Floer
homology and the Rasmussen invariant.

\subsection{14.4 Freedman and Topologically Slice
Knots}\label{ch1986michaelfreedman:freedman-and-topologically-slice-knots}

Freedman's disk theorem shows that every knot with Alexander polynomial
1 is topologically slice. This includes many knots that are not known to
be smoothly slice, such as certain pretzel knots and satellite knots.

\subsection{14.5 The Slice-Ribbon
Conjecture}\label{ch1986michaelfreedman:the-slice-ribbon-conjecture}

The \textbf{slice-ribbon conjecture} (Fox, 1960s) asks whether every
slice knot is a ribbon knot --- a knot that bounds a disk with only
ribbon singularities. This conjecture remains open and is one of the
central problems in knot theory.

\begin{center}\rule{0.5\linewidth}{0.5pt}\end{center}

\section{15. The Schoenflies
Conjecture}\label{ch1986michaelfreedman:the-schoenflies-conjecture}

\subsection{15.1 Statement}\label{ch1986michaelfreedman:statement-1}

The \textbf{Schoenflies conjecture} states that every embedded
\((n-1)\)-sphere in \(S^n\) bounds an \(n\)-ball on each side. It is a
fundamental result in topology that was proved in dimensions
\(n \geq 5\) by Brown and Mazur in the 1960s and is trivial in dimension
2.

\subsection{15.2 Freedman's Result}\label{ch1986michaelfreedman:freedmans-result}

Freedman proved the topological Schoenflies conjecture in dimension 4:
every topologically embedded 3-sphere in \(S^4\) bounds a topological
4-ball on each side. This result is closely related to the 4-dimensional
Poincare conjecture --- in fact, Freedman's proof of the Poincare
conjecture uses the Schoenflies theorem as a building block to construct
the required homeomorphism.

\subsection{15.3 Smooth Version}\label{ch1986michaelfreedman:smooth-version}

As with the h-cobordism theorem, the smooth version of the Schoenflies
conjecture in dimension 4 is false. There exist smoothly embedded
3-spheres in \(S^4\) that do not bound a smooth 4-ball on both sides.
This is another manifestation of the difference between the topological
and smooth categories in dimension 4.

\subsection{15.4 Applications}\label{ch1986michaelfreedman:applications}

The Schoenflies theorem is used in the classification of four-manifolds
to identify standard neighborhoods of embedded 3-spheres. It is
essential for constructing the topological handlebody decomposition that
Freedman uses to build four-manifolds with prescribed intersection
forms.

\subsection{15.5 Relationship to the Annulus
Theorem}\label{ch1986michaelfreedman:relationship-to-the-annulus-theorem}

The Schoenflies theorem is closely related to the annulus theorem, and
Freedman's proof of both theorems uses similar techniques based on
Casson handles and controlled topology.

\begin{center}\rule{0.5\linewidth}{0.5pt}\end{center}

\section{16. The Annulus Theorem}\label{ch1986michaelfreedman:the-annulus-theorem}

\subsection{16.1 Statement}\label{ch1986michaelfreedman:statement-2}

The \textbf{annulus theorem} states that if \(h: S_1 \to S_2\) is a
homeomorphism between two disjoint codimension-1 spheres embedded in a
manifold, and the region between them is homeomorphic to
\(S^{n-1} \times [0,1]\), then \(h\) extends to a homeomorphism of the
entire region between them. In dimension 4, this concerns embeddings of
3-spheres.

\subsection{16.2 Role in the
Classification}\label{ch1986michaelfreedman:role-in-the-classification}

The annulus theorem is a necessary ingredient in the classification of
four-manifolds, particularly in the construction of handle
decompositions and the proof that certain topological constructions are
well-defined. Freedman proved the annulus theorem in dimension 4 as part
of his overall theory.

\subsection{16.3 Relationship to Other
Results}\label{ch1986michaelfreedman:relationship-to-other-results}

The annulus theorem is closely related to the Schoenflies theorem and
the h-cobordism theorem. Together, these three results form the
foundation of the topological handlebody theory in dimension 4. The
Schoenflies theorem gives the ``local'' result (about individual
spheres), while the annulus theorem gives the ``global'' result (about
pairs of spheres).

\subsection{16.4 Historical Context}\label{ch1986michaelfreedman:historical-context-1}

The annulus theorem was a long-standing open problem in topology. In
dimensions \(\geq 5\), it was proved by Kirby and Siebenmann as part of
their work on the classification of topological manifolds. The dimension
4 case -- the last remaining dimension -- was proved by Freedman using
Casson handles.

\begin{center}\rule{0.5\linewidth}{0.5pt}\end{center}

\section{17. The Nielsen Realization
Problem}\label{ch1986michaelfreedman:the-nielsen-realization-problem}

\subsection{17.1 Statement}\label{ch1986michaelfreedman:statement-3}

The \textbf{Nielsen realization problem} (posed by Jakob Nielsen in the
1920s) asks whether every finite subgroup of the mapping class group of
a closed surface can be realized as a group of actual homeomorphisms (or
diffeomorphisms) of the surface. The mapping class group is the group of
isotopy classes of orientation-preserving homeomorphisms, and the
question is whether a finite group of isotopy classes contains a group
of actual homeomorphisms representing those classes.

\subsection{17.2 Freedman's
Contributions}\label{ch1986michaelfreedman:freedmans-contributions}

Freedman contributed cases of the Nielsen realization problem using his
topological methods. His approach used the topological techniques he
developed for the classification of four-manifolds, particularly the
controlled surgery theory and the theory of Casson handles.

\subsection{17.3 Complete Solution By
Kerckhoff}\label{ch1986michaelfreedman:complete-solution-by-kerckhoff}

The Nielsen realization problem was eventually solved in full generality
by Steven Kerckhoff using hyperbolic geometry and Teichmuller theory.
Kerckhoff showed that all finite subgroups of the mapping class group
can be realized by isometries of a hyperbolic metric on the surface.

\subsection{17.4 Significance}\label{ch1986michaelfreedman:significance-3}

The Nielsen realization problem is fundamental to the understanding of
surface topology and Teichmuller theory. It connects the algebraic
structure of the mapping class group with the geometric structure of
surfaces, and it has applications to the topology of three- and
four-manifolds.

\begin{center}\rule{0.5\linewidth}{0.5pt}\end{center}

\section{18. Topological Quantum
Computing}\label{ch1986michaelfreedman:topological-quantum-computing}

\subsection{18.1 Station Q}\label{ch1986michaelfreedman:station-q}

Since 2005, Freedman has led \textbf{Microsoft Station Q} at UC Santa
Barbara, a Microsoft Research laboratory focused on topological quantum
computation. Station Q brings together mathematicians, physicists, and
computer scientists in a unique interdisciplinary environment. It is one
of the world's leading research groups in the field.

\subsection{18.2 The Anyon Model}\label{ch1986michaelfreedman:the-anyon-model}

\textbf{Topological quantum computing} uses anyons --- exotic
quasiparticles that arise in two-dimensional systems --- to encode
quantum information. When anyons are braided around each other in
two-dimensional space, the braiding operation modifies the quantum state
in a way that depends only on the topology of the braid, not on the
precise geometry of the paths. This topological protection is the key to
fault tolerance.

\subsection{18.3 Why Topological Quantum Computing
Matters}\label{ch1986michaelfreedman:why-topological-quantum-computing-matters}

The main advantage of topological quantum computing is \textbf{fault
tolerance}: the quantum information is stored in global topological
degrees of freedom that are insensitive to local perturbations. Unlike
other approaches to quantum computing, which require extensive error
correction, topological quantum computers are naturally resistant to
decoherence. This makes them an attractive approach to building a
scalable quantum computer.

\subsection{18.4 Freedman's
Contributions}\label{ch1986michaelfreedman:freedmans-contributions-1}

Freedman has made foundational contributions to topological quantum
computing: - The classification of anyon models using modular tensor
categories, providing a mathematical classification of possible
topological phases. - The proof that anyon braiding in the Fibonacci
anyon model is computationally universal, meaning any quantum
computation can be performed by braiding anyons. - The relationship
between the Jones polynomial and quantum computation, showing that the
Jones polynomial is BQP-complete. - The realization of topological
phases in condensed matter systems, particularly the fractional quantum
Hall effect.

\subsection{18.5 Current Status}\label{ch1986michaelfreedman:current-status}

Microsoft Station Q continues to work on the realization of topological
quantum computing. The longest-term goal is to create a topological
qubit using a Majorana zero mode in a topological superconductor, which
would be the first practical implementation of a topological quantum
computer.

\begin{center}\rule{0.5\linewidth}{0.5pt}\end{center}

\section{19. Fractional Quantum Hall
Effect}\label{ch1986michaelfreedman:fractional-quantum-hall-effect}

\subsection{19.1 Background}\label{ch1986michaelfreedman:background}

The \textbf{fractional quantum Hall effect} (FQHE) is a quantum
phenomenon observed in two-dimensional electron gases at low
temperatures and strong magnetic fields. Discovered by Tsui, Stormer,
and Gossard in 1982, the FQHE exhibits quasiparticles with fractional
electric charge and exotic braiding statistics.

\subsection{19.2 Anyons in the FQHE}\label{ch1986michaelfreedman:anyons-in-the-fqhe}

The quasiparticles in FQHE systems are believed to be anyons ---
particles whose braiding statistics is neither bosonic nor fermionic.
The well-known example is the Laughlin state at filling fraction
\(\nu = 1/3\), whose quasiparticles have charge \(e/3\) and abelian
anyonic statistics. More exotic states, such as the Moore-Read state at
\(\nu = 5/2\), are believed to support non-abelian anyons.

\subsection{19.3 Freedman's Work}\label{ch1986michaelfreedman:freedmans-work}

Freedman has made significant contributions to understanding FQHE states
in relation to topological quantum computing: - The relationship between
FQHE wavefunctions (Laughlin, Moore-Read, Read-Rezayi states) and
topological quantum field theory. - The classification of FQHE states
using modular tensor categories as the algebraic structure underlying
these states. - The potential use of FQHE states as a physical platform
for topological quantum computing, particularly the \(\nu = 5/2\) state.
- The connection between FQHE wavefunctions and the Jones polynomial.

\subsection{19.4 Experimental Status}\label{ch1986michaelfreedman:experimental-status}

The \(\nu = 5/2\) FQHE state is currently the leading candidate for
realizing non-abelian anyons in the laboratory. If confirmed, this would
be a major step toward topological quantum computing.

\begin{center}\rule{0.5\linewidth}{0.5pt}\end{center}

\section{20. Modular Tensor Categories}\label{ch1986michaelfreedman:modular-tensor-categories}

\subsection{20.1 Definition}\label{ch1986michaelfreedman:definition-3}

\textbf{Modular tensor categories} (MTCs) are braided tensor categories
with a nondegenerate braiding. They are the algebraic structure
underlying: - The representation theory of quantum groups at roots of
unity. - Rational conformal field theory and the Verlinde formula. -
Topological quantum field theory (TQFT) in dimension 3. - Anyon models
in condensed matter physics.

\subsection{20.2 Freedman's Work}\label{ch1986michaelfreedman:freedmans-work-1}

Freedman has made extensive use of MTCs in his work on topological
quantum computing: - The algebraic foundation for topological quantum
computing --- every anyon model is described by a unitary MTC. - A
classification tool for anyon systems: MTCs are believed to classify
topological phases of matter in 2+1 dimensions. - A source of invariants
of three-manifolds and links, via the Reshetikhin--Turaev construction
from the MTC of a quantum group.

\subsection{20.3 The Modular Tensor Category
Classification}\label{ch1986michaelfreedman:the-modular-tensor-category-classification}

Classifications of MTCs are an active area of research. Freedman has
contributed to: - Determining which MTCs can arise from physical systems
and which are ``mathematically possible'' but physically unrealizable. -
The completeness of the classification of MTCs of small rank, up to the
classification of the corresponding conformal field theories. - The
relationship between MTCs and conformal field theory via the modular
functor. - The construction of new MTCs from known ones via operations
like Deligne's tensor product and Drinfeld's center construction.

\subsection{20.4 Key Examples}\label{ch1986michaelfreedman:key-examples-1}

Important MTCs in Freedman's work include: - The Fibonacci category (the
simplest non-abelian MTC, with two simple objects). - The Ising category
(related to the \(\nu = 5/2\) FQHE state). - The categories associated
to quantum \(SU(2)\) at roots of unity, which produce the Jones
polynomial.

\begin{center}\rule{0.5\linewidth}{0.5pt}\end{center}

\section{21. Jones Polynomial and Quantum
Computing}\label{ch1986michaelfreedman:jones-polynomial-and-quantum-computing}

\subsection{21.1 Computational
Complexity}\label{ch1986michaelfreedman:computational-complexity}

Freedman, Larsen, and Wang proved a landmark result in quantum
complexity theory: the value of the Jones polynomial at a root of unity
can be efficiently approximated by a quantum computer, and this
approximation problem is \textbf{BQP-complete} --- meaning it is as hard
as any problem that can be solved by a quantum computer.

\subsection{21.2 The Jones Polynomial as a Quantum
Invariant}\label{ch1986michaelfreedman:the-jones-polynomial-as-a-quantum-invariant}

The Jones polynomial \(V_K(q)\) is evaluated at \(q = e^{2\pi i / r}\)
where \(r\) is an integer. The computational problem is: given a knot
\(K\) described by a braid word, approximate \(V_K(q)\) to within a
specified precision. Freedman showed that this is exactly the
computation performed by braiding anyons in a topological quantum
computer.

\subsection{21.3 Implications}\label{ch1986michaelfreedman:implications}

This result has far-reaching implications: - \textbf{Knot theory meets
complexity theory:} Evaluating topological invariants is inherently
quantum mechanical, not classical. - \textbf{Topological quantum
computing:} Anyon braiding, which computes the Jones polynomial, can
perform universal quantum computation for appropriate anyon models. -
\textbf{Quantum algorithms:} The Jones polynomial is one of the few
problems with a provable exponential speedup on a quantum computer
compared to classical computers. - \textbf{Connection to physics:} The
result provides a physical interpretation of quantum computation in
terms of topological phases of matter.

\subsection{21.4 BQP-Completeness}\label{ch1986michaelfreedman:bqp-completeness}

The BQP-completeness of the Jones polynomial means that any quantum
computation can be reduced to the evaluation of a topological invariant.
This suggests that topology and quantum computation are closely connected
at a fundamental level.

\begin{center}\rule{0.5\linewidth}{0.5pt}\end{center}

\section{22. The Freedman--Donaldson
Complementarity}\label{ch1986michaelfreedman:the-freedmandonaldson-complementarity}

\subsection{22.1 Topological vs Smooth}\label{ch1986michaelfreedman:topological-vs-smooth}

Freedman and Donaldson share the 1986 Fields Medal for work that is
closely complementary. Freedman showed that topological four-manifolds
can be classified cleanly using algebraic invariants. Donaldson showed
that smooth four-manifolds are vastly more restricted.

\subsection{22.2 The Combined Picture}\label{ch1986michaelfreedman:the-combined-picture}

Together, their work reveals: - The topological classification is
complete and simple: intersection form + Kirby--Siebenmann invariant. -
The smooth classification is extremely complex: most topological
four-manifolds admit no smooth structure. - Dimension 4 is unique: in no
other dimension does such a gap exist between the topological and smooth
categories.

\subsection{22.3 Ongoing Research}\label{ch1986michaelfreedman:ongoing-research}

The tension between topological and smooth four-manifolds continues to
drive research: - The smooth Poincare conjecture in dimension 4 remains
open. - The existence of exotic smooth structures on \(S^4\) is unknown.
- The classification of smooth four-manifolds is far from complete.

\begin{center}\rule{0.5\linewidth}{0.5pt}\end{center}

\section{23. Three-Manifolds and
Four-Manifolds}\label{ch1986michaelfreedman:three-manifolds-and-four-manifolds}

\subsection{23.1 The Relationship}\label{ch1986michaelfreedman:the-relationship}

Three-manifold topology and four-manifold topology are closely connected.
The boundary of a four-manifold is a three-manifold, and many questions
about four-manifolds reduce to questions about three-manifolds.

\subsection{23.2 Heegaard Splittings}\label{ch1986michaelfreedman:heegaard-splittings}

A \textbf{Heegaard splitting} of a three-manifold decomposes it into two
handlebodies. The study of Heegaard splittings of the boundary
three-manifold is closely related to the handlebody decomposition of the
four-manifold.

\subsection{23.3 Casson's Invariant}\label{ch1986michaelfreedman:cassons-invariant}

Andrew Casson (whose handles are central to Freedman's work) introduced
an invariant of homology 3-spheres. The Casson invariant counts the
number of representations of the fundamental group into \(SU(2)\) and is
related to the topology of four-manifolds bounded by the homology
sphere.

\subsection{23.4 Floer Homology}\label{ch1986michaelfreedman:floer-homology}

Floer homology for three-manifolds, developed by Andreas Floer, is an
infinite-dimensional homology theory that arises from the study of
four-manifold invariants. Instanton Floer homology and Heegaard Floer
homology are powerful tools for studying both three- and four-manifolds.

\begin{center}\rule{0.5\linewidth}{0.5pt}\end{center}

\section{24. Open Problems in Four-Manifold
Topology}\label{ch1986michaelfreedman:open-problems-in-four-manifold-topology}

\subsection{23.1 The Smooth Poincare Conjecture in Dimension
4}\label{ch1986michaelfreedman:the-smooth-poincare-conjecture-in-dimension-4}

Does there exist a smooth homotopy 4-sphere that is not diffeomorphic to
\(S^4\)? This is one of the important open problems in topology.
Freedman's work shows the topological case is true --- every homotopy
4-sphere is homeomorphic to \(S^4\) --- but the smooth case remains
completely open.

\subsection{23.2 The Classification of Smooth
Four-Manifolds}\label{ch1986michaelfreedman:the-classification-of-smooth-four-manifolds}

Unlike the topological case, there is no complete classification of
smooth simply connected four-manifolds. The Seiberg--Witten invariants
(developed after Donaldson's work) provide powerful tools for
distinguishing smooth structures, but the classification is far from
complete. In particular, the geography problem for smooth four-manifolds
--- which pairs of Chern numbers are realized by smooth four-manifolds
--- is only partially solved.

\subsection{\texorpdfstring{23.3 Exotic
\(S^4\)}{23.3 Exotic S\^{}4}}\label{ch1986michaelfreedman:exotic-s4}

It is unknown whether \(S^4\) itself admits any exotic smooth structure.
This is equivalent to the smooth Poincare conjecture in dimension 4 and
is one of the central open problems in topology.

\subsection{23.4 The Slice-Ribbon
Conjecture}\label{ch1986michaelfreedman:the-slice-ribbon-conjecture-1}

This conjecture in knot theory --- that every slice knot is a ribbon
knot --- remains open and is connected to the topology of four-manifolds
via the work of Freedman and others. It represents one of the deep
connections between knot theory and four-manifold topology.

\subsection{23.5 The Triangulation
Conjecture}\label{ch1986michaelfreedman:the-triangulation-conjecture}

It is unknown whether every topological four-manifold admits a
triangulation. This is related to the Kirby--Siebenmann invariant and
the work of Freedman.

\begin{center}\rule{0.5\linewidth}{0.5pt}\end{center}

\section{25. Awards and Recognition}\label{ch1986michaelfreedman:awards-and-recognition}

\textbf{1986:} - Fields Medal (International Mathematical Union) ---
awarded at the ICM in Berkeley for the topological classification of
four-manifolds and the proof of the 4-dimensional Poincare conjecture.
Freedman shared the medal with Donaldson and Faltings.

\textbf{1997:} - Elected to the National Academy of Sciences (USA) ---
one of the highest honors for a scientist in the United States.

\textbf{2003:} - Elected to the American Academy of Arts and Sciences
--- another prestigious academic honor.

\textbf{2005:} - MacArthur Fellowship (the ``mathematician grant'') --- awarded
for his work in topology and its applications to quantum computing.

\textbf{2006:} - National Medal of Science --- the highest scientific
honor in the United States, awarded by the President for his
contributions to topology and the classification of four-manifolds.

\textbf{Additional honors:} - Fellow of the American Mathematical
Society. - Honorary Doctorate, University of Kiel. - Honorary Doctorate,
University of California, Santa Cruz.

\begin{center}\rule{0.5\linewidth}{0.5pt}\end{center}

\section{26. Selected Major
Publications}\label{ch1986michaelfreedman:selected-major-publications}

\begin{enumerate}
\def\labelenumi{\arabic{enumi}.}
\item
 \textbf{``The topology of four-dimensional manifolds.''} \emph{Journal
 of Differential Geometry} 17, 357--453, 1982. --- \emph{The full
 classification of simply connected topological four-manifolds. This is
 the paper that earned Freedman the Fields Medal.}
\item
 \textbf{``The disk theorem for four-dimensional manifolds.''}
 \emph{Proceedings of the International Congress of Mathematicians},
 Warsaw, 1983. --- \emph{The main technical result: Casson handles are
 homeomorphic to standard handles.}
\item
 \textbf{``Topology of 4-manifolds.''} (with F. Quinn) Princeton
 University Press, 1990. --- \emph{The definitive monograph on the
 topological classification of four-manifolds.}
\item
 \textbf{``Topological quantum computation.''} (with M. Larsen, Z.
 Wang) \emph{Communications in Mathematical Physics} 227, 2002. ---
 \emph{The foundations of topological quantum computing, showing that
 anyon braiding can perform universal quantum computation.}
\item
 \textbf{``A new proof of the topological Poincare conjecture in
 dimension 4.''} \emph{Bulletin of the American Mathematical Society}
 11, 1984. --- \emph{The proof of the 4-dimensional Poincare
 conjecture.}
\item
 \textbf{``The work of Casson on four-dimensional manifolds.''} (with
 F. Quinn) \emph{Topology} 23, 1984. --- \emph{The theory of Casson
 handles and their role in four-manifold topology.}
\item
 \textbf{``The classification of four-manifolds.''} \emph{Notices of
 the American Mathematical Society} 34, 1987. --- \emph{A survey of the
 classification results aimed at a broad mathematical audience.}
\item
 \textbf{``Simulation of topological field theories by quantum
 computers.''} (with A. Kitaev, Z. Wang) \emph{Communications in
 Mathematical Physics} 227, 587--603, 2002. --- \emph{The relationship
 between topological quantum field theory and quantum computation,
 showing that TQFTs can be simulated efficiently by quantum computers.}
\end{enumerate}

\begin{center}\rule{0.5\linewidth}{0.5pt}\end{center}

\begin{center}\rule{0.5\linewidth}{0.5pt}\end{center}

\section*{Exercises}
\addcontentsline{toc}{section}{Exercises}

\begin{exercise} Show that the intersection form of $\\mathbb{CP}^2$ is $\\begin{pmatrix} 1 \\end{pmatrix}$ and that of $\\overline{\\mathbb{CP}}^2$ is $\\begin{pmatrix} -1 \\end{pmatrix}$. \end{exercise}

\begin{exercise} Verify that the $E_8$ lattice is positive definite and even, and classify all such forms of rank 8. \end{exercise}

\begin{exercise} Prove that two simply connected 4-manifolds with the same intersection form are homeomorphic (Freedman's theorem). \end{exercise}

\begin{exercise} Construct an exotic $\\mathbb{R}^4$ using Freedman's techniques. \end{exercise}

\begin{center}\rule{0.5\linewidth}{0.5pt}\end{center}

\section*{Summary}
\addcontentsline{toc}{section}{Summary}

Freedman proved that simply connected 4-manifolds are classified up to homeomorphism by their intersection form and Kirby-Siebenmann invariant. This classification, together with Donaldson's differential-geometric results, revealed the extraordinary complexity of smooth 4-manifold topology and led to the discovery of exotic smooth structures on $\mathbb{R}^4$.

\section{27. References and Further
Reading}\label{ch1986michaelfreedman:references-and-further-reading}

\subsection{Official Sources}\label{ch1986michaelfreedman:official-sources}

\begin{itemize}
\tightlist
\item
 IMU Fields Medal 1986 Announcement:
 https://www.mathunion.org/imu-awards/fields-medal/fields-medal-1986
\item
 National Medal of Science citation: National Science Foundation, 2006.
\end{itemize}

\subsection{Key Technical References}\label{ch1986michaelfreedman:key-technical-references}

\begin{itemize}
\tightlist
\item
 Freedman, M. ``The topology of four-dimensional manifolds.''
 \emph{Journal of Differential Geometry} 17, 357--453, 1982.
\item
 Freedman, M. and Quinn, F. \emph{Topology of 4-Manifolds}. Princeton
 Mathematical Series 39, Princeton University Press, 1990.
\item
 Kirby, R. and Siebenmann, L. \emph{Foundational Essays on Topological
 Manifolds, Smoothings, and Triangulations}. Annals of Mathematics
 Studies 88, Princeton University Press, 1977.
\item
 Donaldson, S. and Kronheimer, P. \emph{The Geometry of
 Four-Manifolds}. Oxford University Press, 1990.
\item
 Freedman, M., Larsen, M., and Wang, Z. ``Topological quantum
 computation.'' \emph{Communications in Mathematical Physics} 227,
 587--603, 2002.
\item
 Casson, A. ``Three lectures on new infinite constructions in
 4-dimensional manifolds.'' \emph{Notes from the Orsay Conference on
 Four-Manifolds}, 1974.
\end{itemize}

\subsection{Textbooks and Surveys}\label{ch1986michaelfreedman:textbooks-and-surveys}

\begin{itemize}
\tightlist
\item
 Gompf, R. and Stipsicz, A. \emph{4-Manifolds and Kirby Calculus}.
 Graduate Studies in Mathematics 20, American Mathematical Society,
 1999.
\item
 Kirby, R. \emph{The Topology of 4-Manifolds}. Lecture Notes in
 Mathematics 1374, Springer, 1989.
\item
 Scorpan, A. \emph{The Wild World of 4-Manifolds}. American
 Mathematical Society, 2005.
\item
 Freedman, M. ``The classification of four-manifolds.'' \emph{Notices
 of the American Mathematical Society} 34, 1987.
\item
 Wall, C.T.C. \emph{Surgery on Compact Manifolds}. London Mathematical
 Society Monographs 1, Academic Press, 1970.
\end{itemize}

\subsection{Biography and Exposition}\label{ch1986michaelfreedman:biography-and-exposition}

\begin{itemize}
\tightlist
\item
 Thurston, W. ``The Work of Michael Freedman.'' ICM 1986 Laudatio,
 Berkeley.
\item
 Quinn, F. ``The classification of four-manifolds.'' \emph{Notices of
 the American Mathematical Society} 33, 1986.
\item
 Anon. ``Michael Freedman --- Biography.'' \emph{US National Academy of
 Sciences}, 1997.
\item
 Hatcher, A. ``Four-manifold topology.'' \emph{Mathematical
 Intelligencer} 7, 1985.
\item
 Milnor, J. ``The Poincare conjecture.'' \emph{Collected Papers of John
 Milnor}, 2007.
\end{itemize}

\subsection{Further Reading on Topological Quantum
Computing}\label{ch1986michaelfreedman:further-reading-on-topological-quantum-computing}

\begin{itemize}
\tightlist
\item
 Kitaev, A. ``Fault-tolerant quantum computation by anyons.''
 \emph{Annals of Physics} 303, 2--30, 2003.
\item
 Nayak, C. et al.~``Non-Abelian anyons and topological quantum
 computation.'' \emph{Reviews of Modern Physics} 80, 1083, 2008.
\end{itemize}

\begin{center}\rule{0.5\linewidth}{0.5pt}\end{center}


\begin{center}\rule{0.5\linewidth}{0.5pt}\end{center}

\section*{See Also}
\addcontentsline{toc}{section}{See Also}

This chapter relates to the following work by other Fields Medalists:
\begin{itemize}
  \item \textbf{\hyperref[ch1986simondonaldson:simon-k.-donaldson-fields-medal-1986-a-comprehensive-report]{See Chapter on related work}} --- Donaldson's invariants distinguish exotic structures.
  \item \textbf{\hyperref[ch2006grigoriperelmandeclined:grigori-perelman-fields-medal-2006-declined]{See Chapter on related work}} --- Perelman proved Poincare, related to Freedman's classification.
  \item \textbf{\hyperref[ch2018peterscholze:peter-scholze-fields-medal-2018-a-comprehensive-report-on-his-work-and-impact]{See Chapter on related work}} --- Scholze's work on p-adic geometry relates to Freedman.
\end{itemize}

\emph{(End of report)}
